\section{Formal Definitions}\label{sec:formal_definitions}

\subsection{Codespace Projector}
Given a commuting set $\mathcal G$ of Paulis,
define
\begin{equation}
\upPi_{\mathcal G} :=\prod_{P \in \mathcal G} \frac{\I +P}2
\label{eq:pj_defn}
\end{equation}
as the projector onto their simultaneous +1 eigenspace.
Specifically,
if $\mathcal G$ generates a stabiliser code,
$\upPi_{\mathcal G}$ projects onto its codespace.
In most cases in these appendices,
the representative of a logical operator
appears multiplied by $\upPi_{\mathcal G}$
so the choice of representative makes no difference;
when it is not, we specify the representative.

\subsection{Faults}\label{sec:faults}
We adapt the definitions in Beverland et al.~\cite[\S 3]{Beverland2024}
and more formally define some concepts introduced in \cref{sec:fault_distance}.
Given a stabiliser circuit,
define a set $\mathcal{F}$ of independent events called \emph{faults}.
Each fault $f$ specifies a pair $(P_f, \v\sigma_f)$,
where the \emph{resultant effect} $P_f$ is a Pauli acting on the data qubits,
and the \emph{signature} $\v\sigma_f$ is a binary vector
whose support indicates the detectors the fault flips.
Each fault induces a function $\pi_f(p)$
that outputs its occurrence probability,
and assume $\pi_f(p) =\upTheta(p)$ for noise level $p \ll 1$.

To construct $\mathcal{F}$ from a set of independent error events,
propagate each one to the end of the circuit,
thereby associating it with a pair $(P, \v \sigma)$.
For each distinct pair,
define one fault representing all its associated error events.
Its occurrence probability $\pi_f(p)$ is
the probability an odd number of those events occur.
Coarse-graining from error events to faults
makes the error analysis of circuits much faster.

A \emph{$|\mathcal C|$-fault configuration} (or \emph{configuration} for short)
is a combination $\mathcal C \subseteq \mathcal F$ of faults whose
resultant effect, signature, and occurrence probability is given respectively by
\begin{align}
	P(\mathcal C) &=\prod_{f \in \mathcal C} P_f \\
	\v\sigma(\mathcal C) &=\bigoplus_{f \in \mathcal C} \v\sigma_f \\
	\pi_\mathcal C(p)
	&=`\Big[\prod_{f \in \mathcal C} \frac{\pi_f(p)}{1 -\pi_f(p)}]
	\prod_{f \in \mathcal F} [1 -\pi_f(p)]
	\label{eq:configuration_probability}
\end{align}
where the second product in \cref{eq:configuration_probability}
is independent of $\mathcal C$.